\section{Copies are confined to the leaf diagonal}
\label{sec:copies}

The distance proof used the prescribed anchor representatives. To
bound the number of copies in the unedited matrix $A_h$, we must now
locate every possible copy. Its $64\times64$ prefix will lie in the
anchor groups of a single mode, and its two remaining rows and columns
will then belong to that mode's designated sets.

\subsection{The anchor and the non-anchor traces}

The anchor matrix $S$ has two complementary features. Its lower-right
corner contains every binary word of length five. We will show that
non-anchor rows cannot realize all these words on five non-anchor
columns. This will force more than $32$ positions on at least one axis of any copy to be anchors.
The dense first $32$ rows and columns, contrasted with the sparse
variable-to-anchor incidences and the zero entries between different
modes, will then force the whole copy into one mode.

\begin{lemma}[Properties of the anchor]
\label{lem:anchor-properties}
The rows of $S$ are pairwise distinct, as are its columns. Every column
has at most one zero in the first $32$ rows, and every row has at most
one zero in the first $32$ columns. Each row has at least $58$ ones and
each column has at least $48$ ones. Finally, the restrictions of rows
$33,\ldots,64$ to columns $60,\ldots,64$ are the $32$ distinct binary
words of length five.
\end{lemma}

\begin{proof}
The replacement in the definition of $S$ does not meet the first
$32$ rows or the first $59$ columns. In the first $59$ columns, each
of rows $1,\ldots,59$ has its own distinguishing zero, while rows
$60,\ldots,64$ consist entirely of ones. The latter five rows are
distinct in the replacement block. Thus the rows are distinct, and
each has at least $58$ ones.

Each of the first $59$ columns has exactly one zero, in its own row.
The last five columns are distinct because they are distinct
coordinates on the list of all binary words of length five. Each has
exactly $16$ zeros, all in the replacement block, so it also differs
from every earlier column and has $48$ ones. The remaining assertions
follow directly from the location and definition of that block.
\end{proof}

\begin{lemma}[Incidence with the modes]
\label{lem:mode-incidences}
A variable position on either axis belongs to designated sets of at
most four modes. Consequently, against any collection containing at
most one representative of each anchor group on the opposite axis,
a variable position has at most four entries equal to one.
\end{lemma}

\begin{proof}
Consider a non-leaf row in $R_i^\sigma(p)$. The only possible modes
are $V_i^\sigma$, $V_{i+1}^\sigma$, and $W_{i,d}^\sigma$, where
$d=p\bmod2$; modes whose level is outside $1,\ldots,h$ are omitted.
A leaf row belongs to $V_h^+$, $V_h^-$, and the two modes
$W_{h,d}^+$, $W_{h,d}^-$. There are at most four in either case.

For a non-leaf column in $C_i^\sigma(p)$, the possible modes are
$V_{i+1}^\sigma$, $W_{i+1,0}^\sigma$,
$W_{i+1,1}^\sigma$, and $W_{i,d}^\sigma$, again omitting undefined
levels. A leaf column belongs only to $W_{h,d}^+$ and $W_{h,d}^-$.
For each participating mode, the entries to its anchor groups form
one unit vector. This proves the last assertion as well.
\end{proof}

A set of columns is \emph{shattered} by a set of rows if the
restrictions of those rows to the columns include every binary word
of the corresponding length.

\begin{lemma}[No five non-anchor columns are shattered]
\label{lem:nonanchor-traces}
For any five distinct non-anchor columns of $A_h$, the non-anchor
rows do not realize all $32$ binary words on those columns.
\end{lemma}

\begin{proof}
If one of the five columns is a dummy column, its entry is one in
every non-anchor row, so the conclusion is immediate. Otherwise all
five columns are variable. Partition them according to their classes
$C_i^\sigma$, counting the common leaf class only once.

Within each such class, order the selected columns by node index
(and arbitrarily within a node block). The entry rules for $A_h$
show that the ones of a variable row form a suffix in each class in
which that row is nonzero. A non-leaf row is nonzero in at most two
classes: its own level and sign, and the preceding level with the
same sign. A leaf row is nonzero in at most three classes: the common
leaf class and the two signs at level $h-1$.

A class containing $a$ selected columns has at most $a$ distinct
nonempty suffixes. Across all classes there are therefore at most
five such suffixes. The trace of a variable row is a union of at most
three of them, giving at most
\[
   \sum_{j=0}^3\binom{5}{j}=26
\]
possible traces. The dummy rows contribute at most one additional
trace. Thus at most $27<32$ words occur.
\end{proof}

\subsection{Every anchor copy uses one mode}

\begin{lemma}[Anchor rigidity]
\label{lem:anchor-rigidity}
Every ordered copy of $S$ in $A_h$ uses the row and column anchor
groups of a single mode $t$. On each axis it uses exactly one position
from each of that mode's $64$ groups, in their prescribed order.
\end{lemma}

\begin{proof}
Fix an ordered copy of $S$. Two positions in a single anchor group
have identical entries throughout $A_h$. By
Lemma~\ref{lem:anchor-properties}, the copy therefore uses at most
one position from any one anchor group. The same reasoning shows
that it uses at most one dummy position on each axis.

Suppose that at most $32$ selected rows and at most $32$ selected
columns are anchors. Since all anchors precede all non-anchor
positions, the last $32$ selected rows and the last five selected
columns are non-anchor positions. Their entries give all $32$ binary
words by Lemma~\ref{lem:anchor-properties}, contrary to
Lemma~\ref{lem:nonanchor-traces}. Hence more than $32$ selected
positions are anchors on at least one axis.

Assume first that more than $32$ selected rows are anchors. In
particular the first $32$ selected rows are anchors. No selected
column can be variable: Lemma~\ref{lem:mode-incidences} would give
at most four ones in these rows, whereas each column of $S$ has at
least $31$ ones in its first $32$ rows. There are consequently at
least $63$ selected anchor columns, since at most one selected
column is dummy.

Any two selected anchor columns have a common one among the first
$32$ selected rows: each has at most one zero there. They must
therefore belong to the same mode, because entries between anchor
groups of different modes are zero. Write $t$ for this common mode.
Every selected anchor row also belongs to $t$. Indeed an anchor row
of another mode would be zero on the at least $63$ selected anchor
columns and would have at most one one in the entire copy,
contradicting Lemma~\ref{lem:anchor-properties}.

A selected dummy column is now impossible: its entries to the
selected anchor rows of mode $t$ contain at most one one, because
its entries to that mode form either a unit vector or the zero
vector. This again contradicts the density of the first $32$ rows
of $S$. Thus all selected columns are anchors of mode $t$. Any
non-anchor row has at most one one against those columns, so every
selected row must also be an anchor of $t$.

If instead more than $32$ selected columns are anchors, the same
argument applies with the axes interchanged. All properties used
after the preceding case distinction are symmetric in the axes:
entries between different anchor modes are zero, the first $32$
rows and columns of $S$ each contain at most one zero in the relevant
restriction, all its rows and columns have more than one one, and
Lemma~\ref{lem:mode-incidences} holds on both axes.

There are $64$ selected positions on each axis and only $64$ anchor
groups in mode $t$. No group is repeated, so every group is used.
Their order in the construction is their order in the copy.
\end{proof}

\begin{corollary}[The two body roles]
\label{cor:body-roles}
In every ordered copy of $H$ in $A_h$, the first $64$ rows and
columns are anchor representatives of one mode $t$. Its two body
rows belong respectively to $\mathsf R_1(t),\mathsf R_2(t)$, and its
two body columns belong respectively to
$\mathsf C_1(t),\mathsf C_2(t)$.
\end{corollary}

\begin{proof}
Apply Lemma~\ref{lem:anchor-rigidity} to the copy's prefix $S$.
The two body rows must have entries $e_1^{\mathsf T}$ and
$e_2^{\mathsf T}$ against these anchor columns. An anchor row of
another mode gives the zero vector; an anchor row of mode $t$ gives
a row of $S$, which has more than one one. Thus neither body row is
an anchor. By the definition of the entries between non-anchor
positions and anchor groups, the vector $e_j^{\mathsf T}$ occurs
exactly on $\mathsf R_j(t)$, for $j=1,2$. The two designated sets are
disjoint. The column argument is identical.
\end{proof}
